\chapter{Beweistexte}
tl;dr:
Einige Texte klingen sehr stark nach PSA, weil wir PSA gewohnt sind.
Wir gehen die Texte durch und sehen, dass PSA nicht die beste Interpretation ist; gleichzeitig gebe ich jeweils eine alternative Interpretation und zeige, warum sie legitim ist.

\section{Jesaja 52-53 - Partizipation im Leiden des Christus}
tl;dr:
Der leidende Gottesknecht ist kein Opfer.
Er leidet nicht stellvertretend, sondern repräsentativ.
Es gibt kein PSA in \bibverse{Isaiah}{52-53}{}
\subsection{Rezeptionshistorie}
tl;dr:
\bibverse{Isaiah}{52-53}{} wird in der Zeit des zweiten Tempels und in der frühen Kirche konsistent als repräsentatives
Leiden einer Messianischen Figur verstanden.
Es gibt keine frühe Interpretation dieses Textes in einem stellvertretenden, juristischen Rahmen.
Das schließt die Verwendungen dieses Textes im NT ein.

\section{Galater 3:13 - Auslöschung des Bundesfluchs}
- Fluch des Gesetzes ist genau das - Fluch des Gesetzes

TODO - ich bin noch unsicher, wie genau der Bundesfluch zu verstehen ist und wie exakt Jesus damit interagiert

\section{Römer 3:25 - Kein Opfer in Sicht}
tl;dr:
ἱλαστήριον (hilasterion) sollte in diesem Vers nicht als Opfer verstanden werden.
Die Optionen sind (a) ἱλαστήριον = der Deckel der Bundeslade, und zwar \textit{als Ort der Begegnung mit Gott}\Footcite{jesus-as-the-mercy-seat}
(b) ἱλαστήριον als Versöhnungsgeschenk aus der römisch-griechischen Kultur\Footcite{lamb-of-the-free}.

TODO: welche dieser Varianten ich tatsächlich bevorzuge ist noch unklar; ich muss zuerst Bailey vollständig lesen und verstehen.

\section{Römer 8:3 - Gottes Urteil über die Sünde}
tl;dr:
Es wird nicht `Jesus für unsere Sünde' verurteilt, sondern `die Sünde im Fleisch'.

TODO: Die Optionen sind mMn (a) Bait-and-Switch Modell wie in Chrysostomus (b) Jesu Wirken ermöglicht uns, nicht weiter im Fleisch unter der Sünde, sondern im Geist unter Gott zu leben, und das wird als `Verurteilung' bezeichnet.

\section{2. Korinther 5:21}

